../cictikz/src/cictikz/data/tex/cictikz_preamble.tex

% Reducing the frequency where the work allows it: a pipeline drawn as
% four registers with a small and a big combinational cloud between
% them. The small cloud settles quickly and its registers run on the
% fast ClkA; the big cloud gets the slow ClkB, and f drops out of the
% dynamic power for everything it clocks. Redrawn from the hand sketch.

\begin{circuitikz}[thick]
  ../cictikz/src/cictikz/data/tex/cictikz_lib.tex

  \newcommand{\dffbox}{
    \draw (0,0) rectangle (1.8,2.4);
    \node[anchor=west] at (0.05,1.95) {$D$};
    \node[anchor=east] at (1.75,1.95) {$Q$};
    \node[anchor=east] at (1.75,0.55) {$\overline{Q}$};
    \draw (1.8,0.55) -- (2.3,0.55);
    \draw (0.5,0) -- (0.8,0.25) -- (1.1,0);
  }

  % four registers at x = 0, 5.4, 10.8, 14.4
  \foreach \x/\c in {0/ClkA, 5.4/ClkB, 10.8/ClkB, 14.4/ClkA} {
    \begin{scope}[shift={(\x,0)}] \dffbox \end{scope}
    \node[anchor=north] at ({\x+0.8},-0.3) {\c};
    \draw ({\x+0.8},0) -- ({\x+0.8},-0.25);
  }

  % the clouds between the first three registers
  \draw[blue] plot [smooth cycle, tension=1]
    coordinates {(2.2,1.2) (2.5,2.4) (3.5,3.0) (4.6,2.5) (4.9,1.3)
                 (4.4,0.2) (3.3,-0.2) (2.5,0.3)};
  \node[armygreen] at (3.55,1.4) {Small};
  \draw[blue] plot [smooth cycle, tension=1]
    coordinates {(7.6,1.2) (7.9,2.5) (9.0,3.1) (10.2,2.6) (10.5,1.3)
                 (10.0,0.1) (8.9,-0.3) (7.9,0.2)};
  \node[red] at (9.05,1.4) {Big};

  % data path through everything
  \draw (-1.0,1.2) -- (0,1.2);
  \draw (1.8,1.2) -- (2.4,1.2);
  \draw (4.7,1.2) -- (5.4,1.2);
  \draw (7.2,1.2) -- (7.8,1.2);
  \draw (10.3,1.2) -- (10.8,1.2);
  \draw (12.6,1.2) -- (14.4,1.2);
  \draw (16.2,1.2) -- (17.2,1.2);

  \node[anchor=north] at (8.1,-1.3) {$f_{ClkB} < f_{ClkA}$};

\end{circuitikz}
\end{document}
